% Préambule commun des schémas TikZ — charte Liouaï
% Chaque schéma commence par % Préambule commun des schémas TikZ — charte Liouaï
% Chaque schéma commence par % Préambule commun des schémas TikZ — charte Liouaï
% Chaque schéma commence par \input{_preambule} puis \begin{document}.
\documentclass[tikz,border=6pt]{standalone}
\usepackage[T1]{fontenc}
\usepackage[utf8]{inputenc}
\usepackage[scaled=0.95]{helvet}
\renewcommand{\familydefault}{\sfdefault}
\usepackage{amsmath,amssymb}
\usepackage[european]{circuitikz}
\usetikzlibrary{arrows.meta,positioning,calc,angles,quotes,decorations.pathreplacing,backgrounds}

% Couleurs du logo Liouaï Digital Education
\definecolor{navy}{HTML}{243743}
\definecolor{vert}{HTML}{28B78D}
\definecolor{vertfonce}{HTML}{167A5C}
\definecolor{vertclair}{HTML}{EAF7F2}
\definecolor{gris}{HTML}{A1A9AE}
\definecolor{gristexte}{HTML}{5C6870}
\definecolor{fond}{HTML}{F3F5F6}
\definecolor{jaune}{HTML}{E5C800}
\definecolor{jauneclair}{HTML}{FFFBD6}

\tikzset{
  every picture/.style={line cap=round, line join=round, color=navy, text=navy},
  fleche/.style={-{Stealth[length=2.6mm,width=2mm]}, line width=1.1pt},
  vecteur/.style={fleche, line width=1.4pt},
  titre/.style={font=\bfseries\small, text=navy},
  note/.style={font=\footnotesize, text=gristexte},
  boite/.style={draw=navy, line width=0.9pt, rounded corners=5pt, fill=fond, align=center, inner sep=6pt},
}
 puis \begin{document}.
\documentclass[tikz,border=6pt]{standalone}
\usepackage[T1]{fontenc}
\usepackage[utf8]{inputenc}
\usepackage[scaled=0.95]{helvet}
\renewcommand{\familydefault}{\sfdefault}
\usepackage{amsmath,amssymb}
\usepackage[european]{circuitikz}
\usetikzlibrary{arrows.meta,positioning,calc,angles,quotes,decorations.pathreplacing,backgrounds}

% Couleurs du logo Liouaï Digital Education
\definecolor{navy}{HTML}{243743}
\definecolor{vert}{HTML}{28B78D}
\definecolor{vertfonce}{HTML}{167A5C}
\definecolor{vertclair}{HTML}{EAF7F2}
\definecolor{gris}{HTML}{A1A9AE}
\definecolor{gristexte}{HTML}{5C6870}
\definecolor{fond}{HTML}{F3F5F6}
\definecolor{jaune}{HTML}{E5C800}
\definecolor{jauneclair}{HTML}{FFFBD6}

\tikzset{
  every picture/.style={line cap=round, line join=round, color=navy, text=navy},
  fleche/.style={-{Stealth[length=2.6mm,width=2mm]}, line width=1.1pt},
  vecteur/.style={fleche, line width=1.4pt},
  titre/.style={font=\bfseries\small, text=navy},
  note/.style={font=\footnotesize, text=gristexte},
  boite/.style={draw=navy, line width=0.9pt, rounded corners=5pt, fill=fond, align=center, inner sep=6pt},
}
 puis \begin{document}.
\documentclass[tikz,border=6pt]{standalone}
\usepackage[T1]{fontenc}
\usepackage[utf8]{inputenc}
\usepackage[scaled=0.95]{helvet}
\renewcommand{\familydefault}{\sfdefault}
\usepackage{amsmath,amssymb}
\usepackage[european]{circuitikz}
\usetikzlibrary{arrows.meta,positioning,calc,angles,quotes,decorations.pathreplacing,backgrounds}

% Couleurs du logo Liouaï Digital Education
\definecolor{navy}{HTML}{243743}
\definecolor{vert}{HTML}{28B78D}
\definecolor{vertfonce}{HTML}{167A5C}
\definecolor{vertclair}{HTML}{EAF7F2}
\definecolor{gris}{HTML}{A1A9AE}
\definecolor{gristexte}{HTML}{5C6870}
\definecolor{fond}{HTML}{F3F5F6}
\definecolor{jaune}{HTML}{E5C800}
\definecolor{jauneclair}{HTML}{FFFBD6}

\tikzset{
  every picture/.style={line cap=round, line join=round, color=navy, text=navy},
  fleche/.style={-{Stealth[length=2.6mm,width=2mm]}, line width=1.1pt},
  vecteur/.style={fleche, line width=1.4pt},
  titre/.style={font=\bfseries\small, text=navy},
  note/.style={font=\footnotesize, text=gristexte},
  boite/.style={draw=navy, line width=0.9pt, rounded corners=5pt, fill=fond, align=center, inner sep=6pt},
}

\begin{document}
\begin{tikzpicture}[scale=1.15]
\newcommand{\riviere}{%
  \fill[vertclair] (-1.9,0) rectangle (2.4,2.6);
  \draw[navy, line width=1pt] (-1.9,0) -- (2.4,0) node[note, right] {rive sud};
  \draw[navy, line width=1pt] (-1.9,2.6) -- (2.4,2.6) node[note, right] {rive nord};
  \foreach \y in {0.6, 1.9} \draw[fleche, gris, line width=0.8pt] (-1.7,\y) -- (-1.0,\y);
  \node[note] at (-1.35, 2.25) {courant};}
% --- Cas 1 : pointer droit vers le nord
\begin{scope}
\riviere
\node[titre] at (0.25, 3.15) {Pointer droit vers le nord};
\coordinate (O) at (0,0.25);
\draw[vecteur, vertfonce] (O) -- node[left] {$\mathbf{v}_{\text{propre}}$} ($(O)+(0,2)$);
\draw[vecteur, gris] ($(O)+(0,2)$) -- node[above] {$\mathbf{c}$} ($(O)+(1,2)$);
\draw[vecteur, navy, line width=1.8pt] (O) -- node[right, pos=0.45] {$\mathbf{v}_{\text{réelle}}$} ($(O)+(1,2)$);
\node[note, align=center] at (0.25, -0.55) {le kayak dérive vers l'est};
\end{scope}
% --- Cas 2 : viser l'amont
\begin{scope}[xshift=6.2cm]
\riviere
\node[titre] at (0.25, 3.15) {Viser l'amont};
\coordinate (O) at (0,0.25);
\coordinate (P) at ($(O)+(-1,1.732)$);
\coordinate (N) at ($(O)+(0,1.732)$);
\draw[gris, densely dashed] (O) -- ($(O)+(0,2.2)$);
\draw[vecteur, vertfonce] (O) -- node[left, pos=0.6] {$\mathbf{v}_{\text{propre}}$} (P);
\draw[vecteur, gris] (P) -- node[above] {$\mathbf{c}$} (N);
\draw[vecteur, navy, line width=1.8pt] (O) -- node[right, pos=0.5] {$\mathbf{v}_{\text{réelle}}$} (N);
\coordinate (H) at ($(O)+(0,1)$);
\pic[draw=vertfonce, line width=0.8pt, angle radius=0.75cm, "$\theta$", angle eccentricity=1.35] {angle = H--O--P};
\node[note, align=center] at (0.25, -0.55) {il traverse tout droit, plus lentement};
\end{scope}
\end{tikzpicture}
\end{document}
